\subsection{The three-branch tacnode defect budget}

\begin{statement}{P4.48}{Proved}
At a planar three-branch critical-curve point with uniform
generic transverse type $A_3$, the local conductor defect
is at least ten. At an ordinary nonplanar three-branch
point, the defect is at least nine. Equality in the
second assertion forces Hessian orders $(2,2,2)$,
projected colengths $(1,1,1)$, descent colength one,
and a quartic square $(a xy+b xz+c yz)^2$ with at
least two nonzero coefficients. The case $abc\ne0$
has rational normalization by P4.47.
\end{statement}
\begin{proof}
For a planar critical curve the ideal is $P=(u,h)$.
The rank-one nilradical $J=P/I$ of the conductor image
is generated by at most two elements over $\mathcal O_C$.
Its normalized image has at least one generator value
on each of the three branches. Modulo the three branch
maximal ideals, the source supplies a vector space of
dimension at most two, while the normalized image
supplies dimension three. Hence its descent colength
$q$ is at least one. The planar symmetric-square and
different calculation in P4.42 gives local defect
$\ge3\delta+q$. Since a planar three-branch point has
$\delta\ge3$, its defect is at least ten.

For an ordinary nonplanar point use actual coordinates
with $P=(xy,xz,yz)$. The argument in P4.47 gives
$f\in P^2$. On each axis, its normalized conormal
lattice is exactly $tK$, so the projected colength
is exactly one. Every Hessian order is at least two.
Write the quartic form in the basis $(xy,xz,yz)$ as
\[
 A(xy)^2+C(xz)^2+E(yz)^2
      +2B(xy)(xz)+2D(xy)(yz)+2F(xz)(yz).
\]
The rank-one Hessians give
$AC=B^2$, $AE=D^2$, $CE=F^2$.
If a diagonal coefficient vanishes, its adjacent
off-diagonal coefficients vanish and the remaining
two-by-two block has rank at most one. If all three
diagonal coefficients are nonzero, the vanishing null
cubic on the $x$-axis gives $F=BD/A$. Higher-degree
terms cannot alter the coefficient linear in $x$
of that normal cubic. Thus the whole quartic is a
square $(a xy+b xz+c yz)^2$.

If the quartic is zero, all three Hessian orders are
at least three, and the local defect is at least
$2+3+9/2=19/2$. If the square has just one nonzero
coefficient, say $a xy$, the $z$-axis Hessian order
is at least three. In the source-generator basis
$(xy,xz,yz)$, the three leading quotient-value rows
have the forms
\[
                      (0,1,0),\qquad(0,0,1),
                      \qquad(0,*,*).
\]
Their rank is at most two, so $q\ge1$ and again
the defect is at least $19/2$.

If at least two coefficients are nonzero, all three
Hessian orders are two. The leading quotient-value
matrix has rank two. For example, up to rescaling
its rows it is
\[
 \begin{pmatrix}-b&a&0\\-c&0&a\\0&-c&b\end{pmatrix}.
\]
It has zero determinant and rank two whenever at
least two of $a,b,c$ are nonzero. Therefore $q\ge1$,
and the local defect is at least $2+3+3+1=9$.
Equality forces $q=1$ and the displayed Hessian and
projection orders. P4.47 applies when all three
quadratic coefficients are nonzero.
\end{proof}
